
\begin{table}[t]
\centering
\small
\caption{Comparison with published approaches for automatically identifying
security weaknesses in LLM tools and MCP servers. Existing approaches either
target predefined vulnerability classes or specific agent--tool attacks.
\project{} instead derives argument-specific security properties from tool
metadata and uses them to dynamically discover semantic implementation bugs.}
\label{tab:motivation}
\resizebox{\textwidth}{!}{
\begin{tabular}{l|c|c|c|c|c|c}
\toprule
\textbf{System} &
\textbf{Tool Impl.} &
\textbf{Dynamic} &
\textbf{Argument Semantics} &
\textbf{Broad Bug Classes} &
\textbf{Semantic Oracle} &
\textbf{Auto. Harness} \\
\midrule

MCP at First Glance~\cite{mcp-first-glance} &
\cmark &
\xmark &
\xmark &
\partialmark &
\xmark &
N/A \\

MCP-SecLint~\cite{mcp-seclint} &
\cmark &
\xmark &
\xmark &
\xmark &
\xmark &
N/A \\

Cisco MCP Scanner~\cite{cisco2025mcpscanner} &
\partialmark &
\partialmark &
\xmark &
\partialmark &
\xmark &
\cmark \\

MCPSecBench~\cite{yang2025mcpsecbench} &
\partialmark &
\cmark &
\xmark &
\partialmark &
\xmark &
\xmark \\

MCP-Scanner~\cite{mcp-scanner} &
\partialmark &
\cmark &
\xmark &
\partialmark &
\xmark &
\cmark \\

% MCPSafetyScanner~\cite{mcp-safety-audit} &
% \partialmark &
% \cmark &
% \partialmark &
% \partialmark &
% \xmark &
% \cmark \\

Chord~\cite{les-dissonances} &
\xmark &
\cmark &
\partialmark &
\xmark &
\partialmark &
\cmark \\

VIPER-MCP~\cite{viper-mcp} &
\cmark &
\cmark &
\partialmark &
\partialmark &
\xmark &
\cmark \\

MCPZoo~\cite{mcpzoo} &
\partialmark &
\partialmark &
\xmark &
\xmark &
\xmark &
\cmark \\

Description-code inconsistency~\cite{li2026don,shi2026description} &
\cmark &
\partialmark &
\partialmark &
\xmark &
\partialmark &
\cmark \\

\midrule
\textbf{\project{}} &
\cmark &
\cmark &
\cmark &
\cmark &
\cmark &
\cmark \\
\bottomrule
\end{tabular}
}
\end{table}
\section{Motivation and Challenges}
\label{sec:moti}


The weak LLM--orchestrator--tool boundary changes how tool vulnerabilities
should be searched for. If the orchestration layer establishes structural
validity but not semantic safety, and LLM-generated values may remain
attacker-influenced, then a tool argument cannot be trusted merely because it
arrived through the normal agent execution path. Consequently, vulnerability
discovery need not begin by asking whether a particular attack can induce the
LLM to generate a malicious value. Instead, \emph{every schema-valid argument
should be treated as potentially adversarial}, and the tool implementation
should be tested against the full security contract expressed by its
interface. Exploit delivery can then be established separately.

This perspective exposes a limitation of existing approaches. As summarized
in Table~\ref{tab:motivation}, most begin with a known weakness or attack and
search for instances of it. MCP at First Glance~\cite{mcp-first-glance}
applies general-purpose and MCP-specific vulnerability checks;
MCP-SecLint~\cite{mcp-seclint} tracks tool arguments toward predefined
security-sensitive operations; and MCP-Scanner~\cite{mcp-scanner} detects a
catalog of MCP threats such as tool poisoning and prompt injection. Dynamic
approaches improve confidence but retain the same coupling: Chord
~\cite{les-dissonances}, for example, executes real tools but uses an oracle
constructed specifically for cross-tool attacks. These systems are effective
for the threats they encode, but their coverage is fundamentally bounded by
the vulnerabilities anticipated by the analysis. A semantically invalid
filesystem access, network destination, or application-specific value may be
security-critical even when it matches no predefined attack pattern or sink.

LLM tools provide a way to search more aggressively. Their metadata already
describes what each argument represents and therefore exposes security
properties unavailable from its primitive type alone. A \texttt{string}
described as a ``path within the workspace'' implies containment; one
described as a remote URL implies constraints on schemes and destinations.
Rather than starting from a vulnerability and asking which tools match it, we
can start from each tool argument, derive the security properties implied by
its semantics, and systematically challenge those properties. This reverses
the direction of existing analyses and introduces four challenges.

\paragraph{C1: Deriving Semantic Security Properties.}
Security checks cannot be selected from primitive types alone. The same
\texttt{string} may represent a path, URL, command, or identifier, each with
different security requirements. We must therefore infer reusable security
properties from the documented meaning of individual arguments and translate
them into adversarial input classes and observable violations.

\paragraph{C2: Specializing Checks to Individual Tools.}
A generic weakness does not fully specify correct behavior. Path containment
depends on the workspace promised by the tool; network safety depends on its
intended destinations and policy. The challenge is to specialize reusable
security knowledge to the concrete contract expressed by each tool without
requiring a manually written analysis for every implementation.

\paragraph{C3: Constructing Oracles Under Partial Visibility.}
Many semantic violations terminate normally, so crashes provide no useful
signal. Moreover, some tool implementations can be inspected while others are
available only through their external interface. We therefore need semantic
oracles that can establish a violation from execution behavior when visibility
is available and from externally observable behavior when it is not.

\paragraph{C4: Going Beyond a Fixed Vulnerability Corpus.}
CWE provides a useful foundation, but it cannot enumerate arbitrary
tool-specific policies. Restricting analysis to a fixed set of sanitizers
would reproduce the coverage limitation of existing rule-based approaches.
The challenge is to specialize or extend reusable security knowledge from tool
metadata while keeping the resulting checks grounded in explicit interface
semantics rather than unconstrained LLM judgment.

These challenges motivate \project{}: a framework that treats exposed tool
arguments as adversarial, derives argument-specific security properties from
their documented semantics, and automatically constructs the inputs and
oracles needed to test them. We separately evaluate whether discovered bugs
can be reached through normal LLM tool-calling paths.
